\section{Introduction and Proof Outline}
\label{sec:introduction}

\subsection{Notation and main result}

All distances, angles, and areas are Euclidean. A unit equilateral triangle
has side length one; its open version is its interior. We write $[x,y]$
for the closed segment, $\conv K$ for the convex hull, and $\interior K$
and $\partial K$ for the interior and boundary. The relative interior
$\operatorname{relint}K$ is its interior within its affine span. The measure $\Haus$ is
one-dimensional Hausdorff measure, which agrees with length on segments.
We use $[t]_+=\max\{t,0\}$ and
$\operatorname{cross}(x,y)=x_1y_2-x_2y_1$ for vectors in $\mathbb R^2$.
A \emph{trace} on a set $E$ means its intersection with $E$.

Let
\[
 O=(0,0),\qquad
 V_i=\left(\cos\frac{i\pi}{3},\sin\frac{i\pi}{3}\right),\qquad
 H=\conv\{V_0,\ldots,V_5\},
\]
where indices lie in $\mathbb Z/6\mathbb Z$.  Write
\[
 e_{i,i+1}=[V_i,V_{i+1}],\qquad r_i=[O,V_i],\qquad M_i=\frac12V_i,
\]
and put
\[
 S=\partial H\cup\bigcup_{i=0}^5r_i.
\]
We call $S$ the \emph{hexagon skeleton}. For $X\in\mathbb R^2$ and $r\ge0$,
write $\mathbb B(X,r)=\{Y:\|Y-X\|\le r\}$ for the closed disk.
\begin{figure}[H]
\centering
\begin{tikzpicture}[scale=2.05]
  \HexCoords
  \DrawHexagon
  \DrawRays
  \fill[CenterBlue] (O) circle (.025);
  \node[figlabel,below left] at (O) {$O$; center role $U_C$};
  \foreach \i in {0,...,5}{
    \fill (V\i) circle (.022);
    \fill[white,draw=RadialTeal,line width=.5pt]
      ($(O)!.5!(V\i)$) circle (.019);
  }
  \node[figlabel,right] at (V0) {$V_0$; role $U_0$};
  \node[figlabel,above right] at (V1) {$V_1$; role $U_1$};
  \node[figlabel,above left] at (V2) {$V_2$; role $U_2$};
  \node[figlabel,left] at (V3) {$V_3$; role $U_3$};
  \node[figlabel,below left] at (V4) {$V_4$; role $U_4$};
  \node[figlabel,below right] at (V5) {$V_5$; role $U_5$};
  \node[figlabel,above] at ($(O)!.5!(V0)$) {$M_0$};
  \node[figlabel,above right] at ($(O)!.68!(V1)$) {$r_1$};
  \node[figlabel,above right] at ($(V0)!.53!(V1)$) {$e_{0,1}$};
\end{tikzpicture}

\caption{The center, vertices, radial segments, and radial midpoints of $H$.}
\label{fig:geometry-roles}
\end{figure}

For a nonempty compact set $K\subset\mathbb R^2$, define its
\emph{equilateral enclosure side} by
\begin{equation}
 \Lambda(K):=\inf\{s>0:K\subset T\text{ for a closed equilateral
 triangle }T\text{ of side }s\}.
 \label{eq:enclosure-definition}
\end{equation}
Then
\begin{equation}
 \Haus(\partial H)=6,\qquad
 \Haus\left(\bigcup_{i=0}^5r_i\right)=6,\qquad \Haus(S)=12.
\label{eq:target-lengths}
\end{equation}

We prove that seven open unit equilateral triangles cannot cover $H$.
Equivalently, seven closed unit equilateral triangles cannot cover $LH$
for any $L>1$. The formal statements and their proofs conclude
Section~\ref{sec:assembly}.

The argument first develops the geometry of the C triangle and V triangles.
It then uses trace length (Section~\ref{sec:trace-length}), area loss
(Section~\ref{sec:area-loss}), selected-gap and supported-midpoint enclosures
(Section~\ref{sec:finite-enclosure}), and a zero-gap nine-point enclosure
(Section~\ref{sec:strategy4-reader}). Supporting calculations are included
with their statements. The exact disk-tangent certificate is proved within
the zero-gap argument; its authenticated machine data and verifiers accompany
the manuscript.
